\documentclass[11pt,a4paper]{article}

\usepackage[margin=2.5cm]{geometry}
\usepackage{amsmath,amssymb,amsfonts}
\usepackage{booktabs}
\usepackage{array}
\usepackage{longtable}
\usepackage{hyperref}
\usepackage{xcolor}

\title{BeamLib Theory Manual \\ Chapter 1: Euler--Bernoulli 2D Beam Element}
\author{BeamLib Development Team}
\date{Phase 1, Batch 2A}

\newcommand{\dd}{\mathrm{d}}
\newcommand{\bmath}[1]{\boldsymbol{#1}}

\begin{document}
\maketitle

\section{Scope and Assumptions}

This chapter derives the 2-node planar Euler--Bernoulli (EB) beam element used in BeamLib for slender beams ($L/h > 20$). The element operates in the global $xz$-plane.

\paragraph{Kinematic assumptions.}
\begin{enumerate}
    \item Plane sections remain plane after deformation.
    \item Cross-section normals remain normal to the deformed beam axis (no shear deformation; this is the defining EB assumption).
    \item Displacements and rotations are infinitesimal (linear theory).
    \item Material is linear elastic, isotropic, homogeneous within the section.
\end{enumerate}

The element is restricted to in-plane bending about the $y$-axis with axial extension along the local $x$-axis. Torsion and out-of-plane bending are not modeled here (they are 3D effects, handled in EB3D).

\section{Sign and DOF Conventions}

\subsection{Coordinate system}

We use a right-handed Cartesian frame with axes $(x, y, z)$ satisfying $\hat{x} \times \hat{y} = \hat{z}$. The 2D problem lives in the $xz$-plane:
\begin{itemize}
    \item $x$: in-plane axial reference direction.
    \item $z$: in-plane transverse direction.
    \item $y$: out-of-plane axis, oriented so that $(\hat{x}, \hat{y}, \hat{z})$ form a right-handed frame ($+\hat{y}$ points \emph{out} of the page in the standard view with $+\hat{x}$ rightward and $+\hat{z}$ upward).
\end{itemize}

\textbf{Important.} Positive $\theta_y$ in BeamLib is defined by the structural-mechanics convention $\theta_y = \mathrm{d}w/\mathrm{d}x$ (see \S\ref{sec:thetay-justification}). This is \emph{opposite} in sign to the strict right-hand-rule rotation about $+\hat{y}$. Throughout this manual ``$\theta_y$'' means the BeamLib structural convention unless explicitly noted.

\subsection{Per-node DOFs}

At every node $i$ we store three DOFs in the order
\[
\bmath{d}_i = \begin{bmatrix} u_{x,i} \\ u_{z,i} \\ \theta_{y,i} \end{bmatrix}.
\]
Positive $u_x$ is in the $+x$ direction (global). Positive $u_z$ is in the $+z$ direction (global). Positive $\theta_y$ follows the BeamLib structural convention $\theta_y = \mathrm{d}w/\mathrm{d}x$ (see \S\ref{sec:thetay-justification}; opposite sign vs.\ right-hand rule about $+\hat{y}$).

\subsection{Element DOF vector}

A 2-node EB2D element has 6 DOFs ordered as
\[
\bmath{d}^e = \begin{bmatrix} \bmath{d}_A \\ \bmath{d}_B \end{bmatrix} = \begin{bmatrix} u_{xA}, & u_{zA}, & \theta_{yA}, & u_{xB}, & u_{zB}, & \theta_{yB} \end{bmatrix}^\mathsf{T}.
\]
Indices in this chapter run $0,\dots,5$ to match the C++ implementation.

\subsection{Sign convention table}

\begin{table}[h!]
\centering
\begin{tabular}{l l l}
\toprule
Quantity & Symbol & Positive direction \\
\midrule
Axial displacement & $u_x$ & along $+\hat{x}$ (global $x$) \\
Transverse displacement & $u_z$ & along $+\hat{z}$ (global $z$) \\
Section rotation & $\theta_y$ & structural convention: $\theta_y = \dd w/\dd x$ (positive slope $\Rightarrow$ positive rotation; opposite sign vs.\ right-hand rule about $+\hat{y}$) \\
Axial force on a node & $F_x$ & along $+\hat{x}$ \\
Transverse force on a node & $F_z$ & along $+\hat{z}$ \\
Nodal moment & $M_y$ & work-conjugate to $\theta_y$ (same convention; positive $M_y$ produces positive $\theta_y$) \\
Slope-rotation relation & $\theta_y(x) = \dd w/\dd x$ & with $w \equiv u_z$ \\
\bottomrule
\end{tabular}
\caption{Sign conventions for the EB2D element.}
\end{table}

\paragraph{Justification of $\theta_y = \dd w / \dd x$ (structural convention).}\label{sec:thetay-justification}
BeamLib adopts the standard structural-mechanics convention: a positive section rotation $\theta_y$ corresponds to a positive slope $\dd w / \dd x$. Geometrically, in the $xz$-plane as drawn ($+\hat{x}$ rightward, $+\hat{z}$ upward), this is a counter-clockwise rotation of the cross-section as the viewer reads the figure.

This is the \emph{opposite} sign from a strict right-hand-rule rotation about $+\hat{y}$. By Rodrigues' formula,
\[
R_y^{\mathrm{RH}}(\theta)\hat{x} = (\cos\theta,\,0,\,-\sin\theta),
\]
so a positive right-hand-rule rotation tilts the cross-section normal toward $-\hat{z}$, not $+\hat{z}$. The BeamLib convention is therefore equivalent to a right-hand-rule rotation about $-\hat{y}$.

We adopt the structural convention because (i) it makes bending kinematics intuitive: positive slope $\Rightarrow$ positive rotation, no mental flip needed when reading a diagram; (ii) it matches the standard Hermite shape function $H_2(\xi) = L(\xi - 2\xi^2 + \xi^3)$, whose derivative at $\xi=0$ is $+1$, enforcing $\theta_{yA} = +\dd u_z/\dd x$ at node $A$; (iii) it matches the MATLAB \texttt{JaphyBeam} reference and most structural-mechanics textbooks. The convention propagates consistently to every $k_l$ sign in \S\ref{sec:stiffness}, the residual definition $\bmath{r}_l = \bmath{k}_l\,\bmath{d}_l$, and the verification signs in \S\ref{sec:verification}.

\textbf{Cross-check.} For a cantilever with downward tip load $F_z < 0$: the beam droops ($u_z(L) < 0$), the slope is negative at the tip ($\dd w/\dd x < 0$), so under the BeamLib convention $\theta_y(L) = \dd w/\dd x < 0$. The test \texttt{test\_eb2d\_patch.cpp} asserts this sign explicitly.

\section{Strain--Displacement Relations}

Let $u(x,z)$, $w(x,z)$ denote total displacement of a material point. EB kinematics:
\[
u(x,z) = u_x(x) - z\,\theta_y(x) = u_x(x) - z\,\frac{\dd w}{\dd x}, \qquad w(x,z) = u_z(x).
\]
The axial strain is
\[
\varepsilon_{xx} = \frac{\partial u}{\partial x} = \frac{\dd u_x}{\dd x} - z\,\frac{\dd^2 w}{\dd x^2}.
\]
The bending curvature is $\kappa = \dd^2 w / \dd x^2 = \dd^2 u_z / \dd x^2$.

The strain energy density per unit length, after integrating over the cross-section, is
\[
\frac{1}{2}\!\int_A E\,\varepsilon_{xx}^2 \,\dd A
= \frac{1}{2}\,E A \!\left(\frac{\dd u_x}{\dd x}\right)^2
+ \frac{1}{2}\,E I_z \!\left(\frac{\dd^2 u_z}{\dd x^2}\right)^2,
\]
where $I_z = \int_A z^2\,\dd A$ is the second moment of area about the $y$-axis (the axis perpendicular to the bending plane, i.e.\ $\hat{y}$). \textbf{In BeamLib we store this as \texttt{SectionProperties::Iz}}, consistent with the 2D xz-plane bending convention.

\paragraph{Section property used.} 2D EB bending in the $xz$-plane uses \texttt{props.Iz}. The 2D--3D cross-validation table in PROJECT\_SPEC \S 13.6 maps 2D \texttt{Iz} $\leftrightarrow$ 3D \texttt{Iy} when comparing the same physical scenario.

\section{Shape Functions}

Use a 2-node element on the natural coordinate $\xi = x / L \in [0, 1]$.

\subsection{Axial (linear)}

\[
N_1^a(\xi) = 1-\xi, \qquad N_2^a(\xi) = \xi.
\]
Axial displacement:
\[
u_x(\xi) = N_1^a(\xi)\,u_{xA} + N_2^a(\xi)\,u_{xB}.
\]

\subsection{Transverse (Hermite cubic)}

Four cubic Hermite functions interpolate $u_z$ and $\theta_y = \dd u_z/\dd x$ at the two ends:
\[
\begin{aligned}
H_1(\xi) &= 1 - 3\xi^2 + 2\xi^3, &\qquad H_2(\xi) &= L\,(\xi - 2\xi^2 + \xi^3), \\
H_3(\xi) &= 3\xi^2 - 2\xi^3,     &\qquad H_4(\xi) &= L\,(-\xi^2 + \xi^3).
\end{aligned}
\]
Transverse displacement:
\[
u_z(\xi) = H_1\,u_{zA} + H_2\,\theta_{yA} + H_3\,u_{zB} + H_4\,\theta_{yB}.
\]
Verification at the endpoints (using $\theta_y = \dd u_z/\dd x = (1/L)\,\dd u_z/\dd \xi$):
\[
u_z(0) = u_{zA},\ \ u_z(1) = u_{zB},\ \ \left.\frac{\dd u_z}{\dd x}\right|_{\xi=0} = \theta_{yA},\ \ \left.\frac{\dd u_z}{\dd x}\right|_{\xi=1} = \theta_{yB}.
\]

\subsection{Curvature operator}

Let $\bmath{d}_b = (u_{zA}, \theta_{yA}, u_{zB}, \theta_{yB})^\mathsf{T}$. Then
\[
\kappa(\xi) = \frac{\dd^2 u_z}{\dd x^2} = \bmath{B}_b(\xi)\,\bmath{d}_b,
\]
with
\[
\bmath{B}_b(\xi) = \frac{1}{L^2}\begin{bmatrix} -6 + 12\xi & L(-4 + 6\xi) & 6 - 12\xi & L(-2 + 6\xi) \end{bmatrix}.
\]

\section{Element Stiffness Matrix Derivation}\label{sec:stiffness}

\subsection{Axial stiffness}

\[
\bmath{k}^a_{2\times 2} = \int_0^L EA \begin{bmatrix} N_1'^2 & N_1' N_2' \\ N_1' N_2' & N_2'^2 \end{bmatrix} \dd x
= \frac{EA}{L}\begin{bmatrix} 1 & -1 \\ -1 & 1 \end{bmatrix}.
\]

\subsection{Bending stiffness}

\[
\bmath{k}^b_{4\times 4} = \int_0^L EI_z \,\bmath{B}_b^\mathsf{T}\bmath{B}_b\,\dd x.
\]
Performing the integral yields the well-known result
\[
\bmath{k}^b = \frac{EI_z}{L^3}
\begin{bmatrix}
12 & 6L & -12 & 6L \\
6L & 4L^2 & -6L & 2L^2 \\
-12 & -6L & 12 & -6L \\
6L & 2L^2 & -6L & 4L^2
\end{bmatrix}.
\]

\subsection{Assembled element stiffness \texorpdfstring{$\bmath{k}_l$}{k\_l} in DOF order \texorpdfstring{$[u_{xA},u_{zA},\theta_{yA},u_{xB},u_{zB},\theta_{yB}]$}{[u\_xA, u\_zA, theta\_yA, u\_xB, u\_zB, theta\_yB]}}

Mapping the axial block to indices $(0, 3)$ and the bending block to indices $(1, 2, 4, 5)$, the nonzero entries of the local $6\times 6$ stiffness $\bmath{k}_l$ are:

\begin{align*}
k_l(0,0) &= k_l(3,3) = \tfrac{EA}{L}, & k_l(0,3) &= k_l(3,0) = -\tfrac{EA}{L}, \\[2pt]
k_l(1,1) &= k_l(4,4) = \tfrac{12 E I_z}{L^3}, & k_l(1,4) &= k_l(4,1) = -\tfrac{12 E I_z}{L^3}, \\[2pt]
k_l(1,2) &= k_l(2,1) = \tfrac{6 E I_z}{L^2}, & k_l(1,5) &= k_l(5,1) = \tfrac{6 E I_z}{L^2}, \\[2pt]
k_l(2,2) &= k_l(5,5) = \tfrac{4 E I_z}{L}, & k_l(2,5) &= k_l(5,2) = \tfrac{2 E I_z}{L}, \\[2pt]
k_l(2,4) &= k_l(4,2) = -\tfrac{6 E I_z}{L^2}, & k_l(4,5) &= k_l(5,4) = -\tfrac{6 E I_z}{L^2}.
\end{align*}
All other entries are zero. The matrix is symmetric.

\subsection{Residual convention}

For this linear element the residual (internal force vector) is
\[
\boxed{\,\bmath{r}^e_l = \bmath{k}_l\,\bmath{d}^e_l.\,}
\]
The Newton residual at the global level is $\bmath{R} = \bmath{R}_{\mathrm{int}} - \bmath{F}_{\mathrm{ext}}$ and the correction equation is $\bmath{K}\,\delta = -\bmath{R}$. With this sign, the convergent state satisfies $\bmath{R}_{\mathrm{int}} = \bmath{F}_{\mathrm{ext}}$.

\section{Consistent Mass Matrix}\label{sec:mass-local}

The consistent mass matrix for EB2D uses the same shape functions:
\[
\bmath{m}^a_{2\times 2} = \int_0^L \rho A
\begin{bmatrix} N_1^a N_1^a & N_1^a N_2^a \\ N_1^a N_2^a & N_2^a N_2^a \end{bmatrix} \dd x
= \frac{\rho A L}{6}\begin{bmatrix} 2 & 1 \\ 1 & 2 \end{bmatrix},
\]
\[
\bmath{m}^b_{4\times 4} = \int_0^L \rho A\,\bmath{H}^\mathsf{T}\bmath{H}\,\dd x
= \frac{\rho A L}{420}
\begin{bmatrix}
156 & 22L & 54 & -13L \\
22L & 4L^2 & 13L & -3L^2 \\
54 & 13L & 156 & -22L \\
-13L & -3L^2 & -22L & 4L^2
\end{bmatrix}.
\]
EB2D mass does \emph{not} include rotary inertia (this is the standard textbook consistent mass for EB and matches PROJECT\_SPEC \S 3.4). Embedded in the $6\times 6$ element DOF order:

\begin{align*}
m_l(0,0) &= m_l(3,3) = \tfrac{\rho A L}{3}, & m_l(0,3) &= m_l(3,0) = \tfrac{\rho A L}{6}, \\
m_l(1,1) &= m_l(4,4) = \tfrac{156\,\rho A L}{420}, & m_l(1,4) &= m_l(4,1) = \tfrac{54\,\rho A L}{420}, \\
m_l(1,2) &= m_l(2,1) = \tfrac{22 L\,\rho A L}{420}, & m_l(1,5) &= m_l(5,1) = -\tfrac{13 L\,\rho A L}{420}, \\
m_l(2,2) &= m_l(5,5) = \tfrac{4 L^2\,\rho A L}{420}, & m_l(2,5) &= m_l(5,2) = -\tfrac{3 L^2\,\rho A L}{420}, \\
m_l(2,4) &= m_l(4,2) = \tfrac{13 L\,\rho A L}{420}, & m_l(4,5) &= m_l(5,4) = -\tfrac{22 L\,\rho A L}{420}.
\end{align*}
This matrix is implemented in Batch 2A as part of \texttt{computeMass}, but the assembly path \texttt{BeamModel::assembleMass} is left for Batch 2B.

\section{Coordinate Transformation}

Each element is integrated in a canonical \emph{local} frame where node $A$ sits at the origin and node $B$ at $(L, 0, 0)$. The transformation $\bmath{T}$ relates global element DOFs to local DOFs.

\subsection{Per-node block}

Let $\Delta\bmath{x} = \bmath{x}_B - \bmath{x}_A$ in global coordinates with components $\Delta x_x$ and $\Delta x_z$ (the $y$ component is zero in the 2D problem). Define
\[
L = \|\Delta\bmath{x}\|,\qquad c = \frac{\Delta x_x}{L},\qquad s = \frac{\Delta x_z}{L}.
\]
The $3\times 3$ per-node block is
\[
\bmath{T}_n =
\begin{bmatrix}
c & -s & 0 \\
s &  c & 0 \\
0 &  0 & 1
\end{bmatrix}.
\]
The element-level $6\times 6$ matrix is block-diagonal: $\bmath{T} = \mathrm{diag}(\bmath{T}_n, \bmath{T}_n)$.

\subsection{Application convention}

By construction $\bmath{T}$ is orthogonal ($\bmath{T}^\mathsf{T}\bmath{T} = \bmath{I}$). BeamLib applies the transformation as
\[
\bmath{d}_l = \bmath{T}\,\bmath{d}_g, \qquad
\bmath{r}_g = \bmath{T}^\mathsf{T}\bmath{r}_l, \qquad
\bmath{k}_g = \bmath{T}^\mathsf{T}\bmath{k}_l\,\bmath{T}.
\]
This convention matches the MATLAB reference (\texttt{JaphyBeam/JaphyEulerBernoulliBeam2D/calcTransformationMatrix.m}) and the user spec verbatim.

\paragraph{Sanity checks.}
\begin{itemize}
    \item Horizontal beam ($\bmath{x}_B - \bmath{x}_A = (L,0,0)$): $c=1$, $s=0$, $\bmath{T}_n = \bmath{I}_3$, so local and global frames coincide.
    \item Orthogonality: $\bmath{T}_n^\mathsf{T}\bmath{T}_n = \mathrm{diag}(c^2+s^2,\,c^2+s^2,\,1) = \bmath{I}_3$ since $c^2+s^2 = 1$.
    \item Energy invariance: $\bmath{d}_g^\mathsf{T}\bmath{k}_g\bmath{d}_g = (\bmath{T}\bmath{d}_g)^\mathsf{T}\bmath{k}_l(\bmath{T}\bmath{d}_g) = \bmath{d}_l^\mathsf{T}\bmath{k}_l\bmath{d}_l$.
    \item The $\theta_y$ DOF is invariant ($T_n(2,2) = 1$) because the rotation axis $\hat{y}$ is preserved by 2D in-plane rotation about $\hat{y}$.
\end{itemize}

\paragraph{Note on Batch 2A scope.}
All Batch 2A tests use beams aligned with the global $x$-axis, so $\bmath{T} = \bmath{I}$ and the transformation is the identity. The non-trivial rotated-beam verification is deferred to Batch 2B (PROJECT\_SPEC \S 2B.5).

\subsection{Reference vector}

\texttt{computeTransformation(xA, xB, refVector)} takes a \texttt{refVector} argument for API uniformity with 3D elements but \textbf{ignores it} in 2D. The 2D rotation about $+\hat{y}$ is fully determined by the in-plane direction $(\Delta x_x, \Delta x_z)$.

\section{Global Assembly}

For each element with connectivity $(A, B)$:
\begin{enumerate}
    \item Build local $\bmath{T}$ from $\bmath{x}_A$, $\bmath{x}_B$.
    \item Gather global element DOFs $\bmath{d}_g^e$ from nodal storage.
    \item $\bmath{d}_l^e = \bmath{T}\,\bmath{d}_g^e$.
    \item Pass $\bmath{x}_A^{loc} = (0,0,0)$, $\bmath{x}_B^{loc} = (L,0,0)$, $\bmath{d}_l^e$ to \texttt{EulerBernoulli2D::computeElement} to obtain $\bmath{r}_l^e$ and $\bmath{k}_l^e$.
    \item $\bmath{r}_g^e = \bmath{T}^\mathsf{T}\bmath{r}_l^e$, $\bmath{k}_g^e = \bmath{T}^\mathsf{T}\bmath{k}_l^e\bmath{T}$.
    \item Scatter $\bmath{r}_g^e$, $\bmath{k}_g^e$ into the free-DOF system using the DofMap. Entries with any fixed DOF on either side are dropped from the free-DOF system (only their RHS contribution remains, see \S\ref{sec:prescribed}).
\end{enumerate}

\section{Prescribed Displacement (Elimination + RHS Correction)}\label{sec:prescribed}

For a fixed DOF whose prescribed value is not zero (the Batch 2B settlement case), the BeamLib implementation uses the elimination plus right-hand-side correction. Partitioning the global system into free ($f$) and restrained ($r$) DOFs gives
\[
\begin{bmatrix} \bmath{K}_{ff} & \bmath{K}_{fr} \\ \bmath{K}_{rf} & \bmath{K}_{rr} \end{bmatrix}
\begin{bmatrix} \bmath{u}_f \\ \bmath{u}_r \end{bmatrix}
=
\begin{bmatrix} \bmath{F}_f^{\mathrm{ext}} \\ \bmath{F}_r^{\mathrm{ext}} + \bmath{R}_r^{\mathrm{reaction}} \end{bmatrix},
\]
where $\bmath{u}_r$ is given. The free-DOF equation is
\[
\bmath{K}_{ff}\,\bmath{u}_f = \bmath{F}_f^{\mathrm{ext}} - \bmath{K}_{fr}\,\bmath{u}_r.
\]
The Phase 1 implementation does \emph{not} explicitly form $\bmath{K}_{fr}$. Instead, the element-local displacement vector $\bmath{d}_l^e = \bmath{T}\bmath{d}_g^e$ is collected with each fixed DOF carrying its prescribed value (\texttt{Node::prescribed}), so the element internal force $\bmath{r}_l^e = \bmath{k}_l^e\bmath{d}_l^e$ already contains the $\bmath{K}_{fr}\bmath{u}_r$ contribution. After mapping back to global and scattering only the free-DOF rows, the residual at a free DOF reads
\[
R_i^{\mathrm{int}} = \sum_{e \ni i} (\bmath{T}^\mathsf{T}\bmath{k}_l^e\bmath{T}\bmath{d}_g^e)_i
                   = \bigl(\bmath{K}_{ff}\bmath{u}_f + \bmath{K}_{fr}\bmath{u}_r\bigr)_i,
\]
and the Newton residual $\bmath{R} = \bmath{R}^{\mathrm{int}} - \bmath{F}^{\mathrm{ext}}$ feeds the correction equation $\bmath{K}_{ff}\delta = -\bmath{R}$, which is exactly the elimination form above. No additional bookkeeping for $\bmath{K}_{fr}$ is required because the contribution is computed implicitly through $\bmath{k}_l^e\bmath{d}_l^e$.

This route assumes one initialization step: any fixed DOF must have \texttt{node.dof[d] = node.prescribed[d]} so that \texttt{scatterFreeDofs} (which only updates free entries) does not contradict the prescribed value. The implementation reads $\bmath{d}_g^e$ via \texttt{Node::totalDof}, which returns \texttt{prescribed[d]} when fixed and \texttt{dof[d]} otherwise; this removes the initialization requirement entirely.

\section{Reaction Recovery}\label{sec:reactions}

BeamLib uses full-system reassembly to recover reactions. The implementation reassembles the element residual into a global-DOF-sized vector (size $N_{\text{nodes}} \cdot n_{\text{DPN}}$, no row/column elimination), then defines
\[
R_d^{\mathrm{reaction}} = R_d^{\mathrm{int,full}} - F_d^{\mathrm{ext}}
\quad\text{for every fixed DOF } d.
\]
The $F_d^{\mathrm{ext}}$ term ensures correctness when an external load is applied at a constrained DOF (e.g.\ a downward concentrated load at a fixed support node): the reaction must absorb both the internal force and the externally applied component. Using $R^{\mathrm{int}}$ alone is incorrect when constrained DOFs carry external loads (Codex Round~3 finding~1). Free-DOF entries of the returned reaction vector are left at zero.

\paragraph{Sign of the recovered reaction.}
The returned reaction is the FE residual at the constrained DOF along the positive DOF direction. For the EB2D structural-convention $\theta_y = \dd w/\dd x$, the relationship to the local internal forces at the support cross section is
\[
R_z^{\mathrm{reaction}} = V_{\mathrm{internal}}^{(\text{left-on-right})}(\text{section}),
\qquad
M_y^{\mathrm{reaction}} = -M_{\mathrm{internal}}^{(=EI w'')}(\text{section}).
\]
The minus sign on $M_y$ comes from the difference between the right-hand-rule moment about $+\hat{y}$ (used in the $\bmath{r}\times\bmath{F}$ static-equilibrium computation) and the BeamLib structural moment work-conjugate to $\theta_y = \dd w/\dd x$ (\S\ref{sec:thetay-justification}); they differ by a sign. The cantilever cross-check used by \texttt{test\_eb2d\_postprocess.cpp} confirms this: with $F_z = -1000$ at the tip, $M_{\mathrm{internal}}(0) = EI w''(0) = -1000$, so the recovered reaction is $M_y^{\mathrm{reaction}} = +1000$.

\section{Internal Forces (End Values)}\label{sec:internal-forces}

Each element exposes a static method \texttt{computeInternalForces} that returns the cross-section forces $N$, $V_z$, $M_y$ at both ends. The starting point is the local element nodal force vector
\[
\bmath{f}_e = \bmath{k}_l\,\bmath{d}_l \in \mathbb{R}^6,
\]
which is the force the two nodes apply to the element (equivalently, the FE residual contribution from this element evaluated at the local end DOFs). The conversion to ``beam diagram'' cross-section forces is
\[
\boxed{\;\begin{aligned}
N_A &= -f_e[0], & V_{z,A} &= +f_e[1], & M_{y,A} &= +f_e[2], \\
N_B &= +f_e[3], & V_{z,B} &= -f_e[4], & M_{y,B} &= -f_e[5].
\end{aligned}\;}
\]
The sign rule: at end $A$, the section is just to the \emph{right} of node $A$ (element interior on the right), so $f_e$ entries at end $A$ already give the ``left material on right material'' force. At end $B$, the section is just to the \emph{left} of node $B$ (element interior on the left), so $f_e$ entries at end $B$ give the opposite (``right material on left material''); the negation restores the cross-section ``left on right'' convention so that $V_z$ and $M_y$ at the two ends of an element share the same sign for constant-shear / linearly-varying-moment loadings. For axial force, the BeamLib convention is ``tension positive,'' which is the negative of the ``left-on-right'' axial component; this introduces the opposite-of-shear sign pattern (flip at $A$, no flip at $B$).

\paragraph{Cantilever cross-check (used in \texttt{test\_eb2d\_postprocess.cpp}).}
A horizontal cantilever fixed at $x=0$ with tip load $F_z = -1000$ at $x=L$ has, by statics,
\[
V_z(x) = +1000, \qquad M_y(x) = +1000\,(L - x), \qquad N(x) = 0.
\]
Therefore $V_{z,A} = V_{z,B} = +1000$ for every element; $M_{y,A}$ of the root element equals $+1000$ and $M_{y,B}$ of the tip element equals $0$. These values match the Codex prompt's acceptance criteria and verify the sign of every term in the table above.

\section{Mass Matrix Assembly}\label{sec:mass-assembly}

The global consistent mass matrix follows the same per-element transformation as the stiffness:
\[
\bmath{M}_g^e = \bmath{T}^\mathsf{T}\,\bmath{M}_l^e\,\bmath{T},
\]
with $\bmath{M}_l^e$ from \S\ref{sec:mass-local}. Scatter into the free-DOF system uses the same DofMap as the stiffness assembly; mass entries that fall on a fixed DOF are dropped. This is appropriate for modal analysis with free-DOF generalized eigenvalue problems (Batch~9).

\section{Formula-to-C++ Index Mapping}

\begin{longtable}{c c l}
\toprule
Formula symbol & C++ index & Meaning \\
\midrule
\endhead
$u_{xA}$ & \texttt{dispVec[0]} & node-A axial displacement (global $x$ in canonical local frame) \\
$u_{zA}$ & \texttt{dispVec[1]} & node-A transverse displacement \\
$\theta_{yA}$ & \texttt{dispVec[2]} & node-A section rotation \\
$u_{xB}$ & \texttt{dispVec[3]} & node-B axial displacement \\
$u_{zB}$ & \texttt{dispVec[4]} & node-B transverse displacement \\
$\theta_{yB}$ & \texttt{dispVec[5]} & node-B section rotation \\
\midrule
$k_l(0,0)$ & \texttt{ke(0,0)} & $EA/L$ \\
$k_l(3,3)$ & \texttt{ke(3,3)} & $EA/L$ \\
$k_l(0,3)$ & \texttt{ke(0,3)} & $-EA/L$ \\
$k_l(1,1)$ & \texttt{ke(1,1)} & $12 E I_z / L^3$ \\
$k_l(1,2)$ & \texttt{ke(1,2)} & $6 E I_z / L^2$ \\
$k_l(1,4)$ & \texttt{ke(1,4)} & $-12 E I_z / L^3$ \\
$k_l(1,5)$ & \texttt{ke(1,5)} & $6 E I_z / L^2$ \\
$k_l(2,2)$ & \texttt{ke(2,2)} & $4 E I_z / L$ \\
$k_l(2,4)$ & \texttt{ke(2,4)} & $-6 E I_z / L^2$ \\
$k_l(2,5)$ & \texttt{ke(2,5)} & $2 E I_z / L$ \\
$k_l(4,4)$ & \texttt{ke(4,4)} & $12 E I_z / L^3$ \\
$k_l(4,5)$ & \texttt{ke(4,5)} & $-6 E I_z / L^2$ \\
$k_l(5,5)$ & \texttt{ke(5,5)} & $4 E I_z / L$ \\
\midrule
$c$ & \texttt{c} & $\Delta x_x / L$ \\
$s$ & \texttt{s} & $\Delta x_z / L$ (note: $z$, not $y$, in the xz-plane) \\
$T_n(0,1)$ & \texttt{T(0,1)} & $-s$ \\
$T_n(1,0)$ & \texttt{T(1,0)} & $+s$ \\
\bottomrule
\caption{Formula symbol to C++ index mapping for EB2D.}
\end{longtable}

\section{Verification: Analytical Solutions}\label{sec:verification}

\subsection{Cantilever tip transverse load}

A cantilever of length $L$ fixed at node 0 with a transverse load $F_z$ at the tip satisfies
\[
u_z(L) = \frac{F_z\,L^3}{3 E I_z}, \qquad \theta_y(L) = \frac{F_z\,L^2}{2 E I_z}.
\]
For $F_z < 0$ (downward), both are negative (tip droops, slope is negative). With $F_z = -F$ ($F > 0$ being the magnitude),
\[
u_z(L) = -\frac{F L^3}{3 E I_z}, \qquad \theta_y(L) = -\frac{F L^2}{2 E I_z}.
\]
Reference numerical values used in \texttt{test\_eb2d\_cantilever.cpp}: $E = 200\!\times\!10^9$, $A = 0.01$, $I_z = 8.333\!\times\!10^{-6}$, $L=1.0$, $F=1000$, giving $u_z(L) = -2\!\times\!10^{-4}$ and $\theta_y(L) = -3\!\times\!10^{-4}$.

\subsection{Cantilever tip end-moment}

A cantilever with an applied end moment $M_y$ at the tip yields a constant curvature $\kappa = M_y / (E I_z)$ throughout the beam and
\[
u_z(L) = \frac{M_y\,L^2}{2 E I_z}, \qquad \theta_y(L) = \frac{M_y\,L}{E I_z}.
\]
The bending moment is constant along the span (no shear, pure bending).

\subsection{Patch test (constant axial strain)}

Apply $u_{xB} = u_0$ at the right end of an arbitrary mesh of EB2D elements aligned with $+\hat{x}$, with the left end pinned axially. Every element should report axial strain $u_0 / L_{\text{total}}$ and zero bending response.

\subsection{Rigid translation}

Setting $u_{xA} = u_{xB}$, $u_{zA} = u_{zB}$, $\theta_{yA} = \theta_{yB} = 0$ should give $\bmath{r}_l = \bmath{k}_l \bmath{d}_l$ where the axial block contributes zero (equal axial DOFs) and bending block also gives zero (constant deflection, zero curvature, zero rotation difference). Hence $\bmath{r}_l = \bmath{0}$.

\section{Verification Tolerance Table}

\begin{table}[h!]
\centering
\begin{tabular}{l l}
\toprule
Check & Tolerance \\
\midrule
Single-element matrix entries ($EA/L$, $12EI/L^3$, $\dots$) & relative $10^{-12}$ \\
Patch test rigid-translation residual & absolute $10^{-12}$ \\
Patch test constant-strain residual & absolute $10^{-10}$ \\
Cantilever tip $u_z$, $\theta_y$ (10 elements) & relative $10^{-9}$ \\
Pure end-moment $u_z$, $\theta_y$ (10 elements) & relative $10^{-9}$ \\
Constant-moment uniformity along span & relative $10^{-9}$ \\
Cantilever reactions, element internal forces (post-process) & relative $10^{-9}$ \\
Simply supported beam midspan deflection (20 elements) & relative $10^{-8}$ \\
Simply supported beam reactions $qL/2$ & relative $10^{-9}$ \\
Rotated beam local displacement vs.\ horizontal reference & relative $10^{-9}$ \\
Settlement: deflection field, reactions (10 elements) & relative $10^{-9}$ \\
Reaction at unloaded free DOF (absolute) & $10^{-9}$ \\
Newton iteration count (linear case) & exactly 1 correction step \\
Final NR residual (linear case) & $\leq 10^{-7}$ absolute (set by NRConfig.tol) \\
\bottomrule
\end{tabular}
\caption{Verification tolerance for EB2D tests.}
\end{table}

\section{Implementation Pitfalls (for the Reviewer)}

\begin{enumerate}
    \item \textbf{Section property choice.} EB2D bending in the $xz$-plane uses \texttt{props.Iz}. Using \texttt{props.Iy} is a silent bug in 2D (the field exists but has no role here).
    \item \textbf{$s$ comes from $\Delta x_z$, not $\Delta x_y$.} The 2D problem lives in the $xz$-plane; $y$ is out-of-plane. Reusing a generic ``$xy$'' rotation formula will give wrong transformations for non-horizontal beams.
    \item \textbf{Residual sign.} The element returns $\bmath{r}_l = \bmath{k}_l \bmath{d}_l$, not $-\bmath{k}_l \bmath{d}_l$. The Newton solver assembles $\bmath{R} = \bmath{R}_{\mathrm{int}} - \bmath{F}_{\mathrm{ext}}$ and solves $\bmath{K}\,\delta = -\bmath{R}$.
    \item \textbf{$\theta_y$ sign under $F_z$.} A downward tip load yields a negative $\theta_y$ at the tip (test \texttt{test\_eb2d\_patch.cpp} explicitly checks this).
    \item \textbf{Bending couplings.} The signs in $k_l$ entries connecting $u_z$ and $\theta_y$ at the same node are positive (e.g.\ $k_l(1,2) = +6EI/L^2$), while couplings between $u_z$ at one node and $\theta_y$ at the other can be positive or negative depending on which side (see the index map above).
\end{enumerate}

\end{document}
